\documentclass[11pt]{article}

\usepackage[margin=1in]{geometry}
\usepackage{array}
\usepackage{booktabs}
\usepackage{amsmath}
\usepackage{hyperref}
\hypersetup{colorlinks=true, linkcolor=blue, urlcolor=blue}
\usepackage{enumitem}

\newcommand{\code}[1]{\texttt{#1}}

\title{Cloud Optimization: TD1\\
\large Local Workloads, Linux Setup, and Benchmarking}
\author{Group 19: Yuba, Amani, and Hichem \\ A4, Cloud Computing \& Cybersecurity (CCC)}
\date{Session: 2026-09-11 \quad Report finalized: 2026-09-17}

\begin{document}
\maketitle

\begin{abstract}
TD1 required no cloud infrastructure. Its purpose was to give us a disciplined, reproducible
local benchmarking methodology (environment characterization, compilation, timing, and a
controlled measurement campaign) before we attempt any cross-machine or cross-VM comparison in
later sessions. In this report we document that methodology as we applied it across our three
laptops, and we present two findings that go beyond what the TD sheet asked for: a cache cliff in
\code{matrix\_task} and a $\sim\!2\times$ runtime effect from CPU power state.
\end{abstract}

\clearpage
\setcounter{tocdepth}{1}
\tableofcontents
\clearpage

\section{Objectives and scope}
Per the TD1 sheet, our learning objectives were to prepare a Linux environment, compile a Go
program to a native executable, characterize our local machines, measure execution time and
memory, run repeated experiments across several workload sizes, and produce a reproducible local
benchmark dataset. TD1 is not graded and has no formal submission, so this report is our own
record, as a group, of what we did and what we found, alongside the raw data in our project
repository.

\section{Environment}

\subsection{Our three machines}
We ran the benchmarks independently on our own laptops so that we would have more than one data
point from the start. Our primary campaign (Sections~\ref{sec:campaign}
to~\ref{sec:finding2}) uses Hichem's machine, \code{ultra9-285h}; Section~\ref{sec:cross-machine}
brings in Amani's and Yuba's machines for comparison.

\begin{center}
\begin{tabular}{llll}
\toprule
& Hichem (\code{ultra9-285h}) & Amani & Yuba \\
\midrule
Distribution & Ubuntu 24.04.4 LTS & (WSL2) & Debian 12 (bookworm) \\
Environment & WSL2 & WSL2 & WSL2 \\
CPU & Core Ultra 9 285H & Core Ultra 7 268V & i7-11800H \\
Cores / threads & 16 physical, 1/core & 8 physical, 1/core & 8 physical, 2/core (SMT) \\
Logical CPUs & 16 & 8 & 16 \\
L3 cache & 24\,MiB & 12\,MiB & 24\,MiB \\
RAM & 19\,GiB & (not recorded) & 15\,GiB \\
Go & go1.22.2 & (see machine file) & go1.19.8 \\
\bottomrule
\end{tabular}
\end{center}

We are all on WSL2, which is a virtual machine on a Windows host rather than bare metal. We
confirmed WSL2 is acceptable for this project, but we still think it is worth declaring
explicitly: it is exactly the kind of environment difference the sheet warns us against ignoring
when comparing results.

\subsection{Why we recorded the environment first}
A runtime is not a property of a program in isolation. It is a property of a program running on a
specific machine, so without the machine recorded, a number cannot be interpreted or reproduced.
This is also why every row of our benchmark CSV (Section~\ref{sec:campaign}) carries a machine
label rather than being pooled anonymously.

\subsection{Core vs.\ logical CPU}
A core is a physical execution unit; a logical CPU is what the OS scheduler can dispatch onto.
They differ only when simultaneous multithreading (SMT, or hyperthreading) lets one core present
as two scheduling targets. Hichem's and Amani's machines both report one thread per core, so core
count equals logical CPU count for us on those two. Yuba's machine is the exception
(Section~\ref{sec:cross-machine}): 8 physical cores but 16 logical CPUs, with SMT on.

\section{Compilation}
We compiled the four reference workloads (\code{cpu\_task}, \code{memory\_task},
\code{matrix\_task}, \code{sort\_task}) with:
\begin{verbatim}
go build -o bin/cpu_task    workloads/cpu_task.go
go build -o bin/memory_task workloads/memory_task.go
go build -o bin/matrix_task workloads/matrix_task.go
go build -o bin/sort_task   workloads/sort_task.go
\end{verbatim}

\code{go build} produces one self-contained native executable per source file ($\sim$2.3\,MB,
\code{ELF 64-bit LSB executable, x86-64, statically linked}; \code{ldd} reports \emph{not a
dynamic executable}). Because Go compiles ahead of time, nothing is translated or interpreted
inside a timed region, so our measurements reflect the work itself, not a runtime. We confirmed the
binaries run standalone, with source and environment stripped:
\begin{verbatim}
$ env -i ./cpu_task --size 100000
{"task":"cpu_task","size":100000,"runtime_s":0.007179392,"primes_found":9592}
\end{verbatim}

\section{The four workloads}
\begin{center}
\begin{tabular}{lllc}
\toprule
Workload & What it does & \code{--size} means & Cost shape \\
\midrule
\code{cpu\_task} & counts primes by trial division & largest integer tested & $n^{1.5}$ \\
\code{memory\_task} & allocates + sweeps an array (5 passes) & megabytes allocated & linear \\
\code{matrix\_task} & naive dense $N \times N$ multiply & matrix dimension $N$ & $N^3$ \\
\code{sort\_task} & generates then sorts an integer array & number of integers & $n\log n$ \\
\bottomrule
\end{tabular}
\end{center}

We found \code{cpu\_task} to be the most CPU-intensive workload: a flat 2\,MB footprint at every
size, and 99\% user time (4.55\,s user versus 0.05\,s system at its largest tested size).
\code{memory\_task} uses the most memory: peak RSS tracks the requested size almost exactly
(1024\,MB requested, 1027\,MB resident). Taking size as a command-line parameter let us produce a
controlled range of work from one binary, which is what makes it possible to measure how cost grows
with input rather than only at one point.

Two implementation details we want to flag for whoever automates this in TD3: \code{memory\_task}
emits a \code{size\_mb} field where the other three emit \code{size}, and the four programs' internal
timers start at different points in their own execution (Section~\ref{sec:timing}).

\section{Timing methodology}
\label{sec:timing}

\subsection{Wall-clock vs.\ CPU time}
Wall-clock time is elapsed real time. CPU time is time a core actually spent on the process, split
into user (its own instructions) and system (kernel work on its behalf); the two diverge whenever
the process is not actively running. \code{memory\_task} is the one workload where we measured
system time exceeding user time: the real cost is the kernel supplying and zeroing a gigabyte of
pages, not the program touching it. It is also the only workload that drops below 100\% CPU,
blocking on page faults.

\subsection{What we decided about \code{runtime\_s}}
We settled two measurement questions ourselves, directly from the sheet's own text and our own
data, without needing to wait on an external answer.

\begin{itemize}[leftmargin=*]
\item \textbf{We used the GNU \code{time} wall clock, not each program's internal JSON figure.}
The sheet defines \code{runtime\_s} as elapsed time from start to end of execution, and the
external wall clock is the only definition consistent across all four workloads. The internal
timers disagree with wall clock by workload: about 1\% for \code{cpu\_task} and
\code{memory\_task}, about 1\% for \code{matrix\_task} (which excludes filling its input
matrices), and 10\% for \code{sort\_task} (which excludes generating its input array). Comparing
internal figures across workloads would have silently compared different things, so we kept both:
the CSV holds wall clock, and we preserved the per-run JSON under \code{data/raw/}.
\item \textbf{We treated one CSV row as one single execution, never a mean.} The sheet's own
Part~6 instructions are explicit: keep the individual measurements, and ``compute the average
runtime only during the analysis phase.'' Our CSV schema
(\code{workload,size,run,runtime\_s,max\_rss\_kb}) has a dedicated \code{run} column precisely so
raw repetitions are never collapsed on write.
\end{itemize}

\subsection{Why two runs of the same program differ}
We identified four causes, largest effect first:
\begin{enumerate}
\item \textbf{CPU power state, $\sim\!2\times$.} The single largest effect we measured in this TD;
see Finding 2, Section~\ref{sec:finding2}.
\item \textbf{The definitional gap between internal and wall-clock timers} (not true variation),
quantified above and largest for \code{sort\_task}.
\item \textbf{Cache effects.} Small between repeats at fixed size, large across sizes; see Finding
1, Section~\ref{sec:finding1}.
\item \textbf{Scheduling and other processes.} Repeats within our campaign spread 1 to 5\%.
\end{enumerate}

\section{Benchmark procedure and local campaign}
\label{sec:campaign}
We wrote a Bash runner, \code{scripts/benchmark.sh}, to automate the campaign: it selects a
workload and size, runs it a configurable number of times, measures wall clock and peak RSS via
\code{/usr/bin/time -v}, appends one CSV row per execution, and reports a clear error (exit
non-zero, no row written) if the binary is missing or a run fails. Every row carries the producing
machine's label, matching a \code{data/machine\_<label>.txt} spec file.

The sheet's suggested sizes all completed in under 0.25\,s on Hichem's CPU, several in under
0.01\,s, indistinguishable from process startup and scheduler noise. We scaled the sizes up until
the largest run took roughly 2\,s:

\begin{center}
\begin{tabular}{lccc}
\toprule
Workload & Sheet suggested & Used here & Largest run \\
\midrule
\code{cpu\_task} & 200k / 500k / 1M & 2.5M / 5M / 10M & 2.12\,s \\
\code{memory\_task} & 128 / 256 / 512\,MB & 256 / 512 / 1024\,MB & 1.13\,s \\
\code{matrix\_task} & 200 / 400 / 600 & 500 / 1000 / 1500 & 4.50\,s \\
\code{sort\_task} & (absent from the sheet's campaign) & 5M / 10M / 20M & 1.87\,s \\
\bottomrule
\end{tabular}
\end{center}

We left \code{--passes} at its default of 5 throughout, since changing it would make sizes
incomparable. The results below are the mean of 3 runs, measured on mains power, across the 36
rows in \code{data/benchmark\_local.csv}.

\begin{center}
\begin{tabular}{lrrrrr}
\toprule
Workload & Size & Mean (s) & Min & Max & Peak RSS \\
\midrule
\code{cpu\_task}    & 2\,500\,000  & 0.283 & 0.28 & 0.29 & 2\,MB \\
\code{cpu\_task}    & 5\,000\,000  & 0.750 & 0.75 & 0.75 & 2\,MB \\
\code{cpu\_task}    & 10\,000\,000 & 2.120 & 2.02 & 2.22 & 2\,MB \\
\code{memory\_task} & 256          & 0.300 & 0.27 & 0.36 & 258\,MB \\
\code{memory\_task} & 512          & 0.547 & 0.49 & 0.63 & 514\,MB \\
\code{memory\_task} & 1024         & 1.127 & 1.05 & 1.28 & 1027\,MB \\
\code{matrix\_task} & 500          & 0.073 & 0.07 & 0.08 & 8\,MB \\
\code{matrix\_task} & 1000         & 0.550 & 0.54 & 0.56 & 25\,MB \\
\code{matrix\_task} & 1500         & 4.503 & 4.13 & 4.73 & 54\,MB \\
\code{sort\_task}   & 5\,000\,000  & 0.423 & 0.41 & 0.43 & 40\,MB \\
\code{sort\_task}   & 10\,000\,000 & 0.863 & 0.85 & 0.89 & 79\,MB \\
\code{sort\_task}   & 20\,000\,000 & 1.870 & 1.83 & 1.91 & 155\,MB \\
\bottomrule
\end{tabular}
\end{center}

\section{Finding 1: a cache cliff in \code{matrix\_task}}
\label{sec:finding1}
We predicted scaling from each workload's cost shape before measuring. Doubling the input does
\emph{not} double the runtime, except for the linear workload:

\begin{center}
\begin{tabular}{lcccc}
\toprule
Workload & Step & Predicted & Observed & \\
\midrule
\code{cpu\_task}    & 5M $\to$ 10M   & $2.83\times$ & $2.83\times$ & matches \\
\code{memory\_task} & 512 $\to$ 1024 & $2\times$    & $2.06\times$ & matches \\
\code{sort\_task}   & 10M $\to$ 20M  & $\sim\!2.1\times$ & $2.17\times$ & matches \\
\code{matrix\_task} & 500 $\to$ 1000 & $8\times$    & $7.53\times$ & matches \\
\code{matrix\_task} & \textbf{1000 $\to$ 1500} & \textbf{$3.38\times$} & \textbf{$8.19\times$} &
\textbf{2.4$\times$ too slow} \\
\bottomrule
\end{tabular}
\end{center}

\code{matrix\_task} holds three dense $N \times N$ float64 matrices, so its working set is
$3N^2 \times 8$ bytes:

\begin{center}
\begin{tabular}{lccc}
\toprule
$N$ & Predicted working set & vs.\ 24\,MiB L3 & Measured RSS \\
\midrule
500  & 6\,MB  & fits             & 8\,MB \\
1000 & 24\,MB & at the boundary  & 25\,MB \\
1500 & 54\,MB & 2.25$\times$ over & 54\,MB \\
\bottomrule
\end{tabular}
\end{center}

Once the working set leaves last-level cache, inner-loop steps become DRAM round-trips, which
explains the jump from $3.38\times$ to $8.19\times$. Our measured RSS matches the predicted
working set at all three sizes, and the $N=1500$ run also has the widest spread in our campaign
(4.13 to 4.73\,s, about 14\%) against about 5\% for \code{cpu\_task}: the signature of a
memory-bound run.

We consider this the finding that matters most for our project. Three of the four workloads can be
ranked across machine types from clock speed alone. \code{matrix\_task} cannot: its cost depends
on cache hierarchy, which is invisible in a vCPU count or a price list. Any scheduler that models
task cost as a single scalar per machine type will mispredict it.

\section{Finding 2: power state moves every runtime by $\sim\!2\times$}
\label{sec:finding2}
We re-measured roughly 90 minutes after our main campaign, on battery instead of mains, and got
$\sim\!2\times$ longer runtimes on an otherwise idle machine.

\begin{center}
\begin{tabular}{lccc}
\toprule
& Mains & Battery & Slowdown \\
\midrule
\code{cpu\_task} 10M                    & 2.120 & 4.257 & $2.01\times$ \\
\code{matrix\_task} 1000 (fits L3)      & 0.550 & 1.367 & $2.49\times$ \\
\code{matrix\_task} 1500 (exceeds L3)   & 4.503 & 5.733 & $1.27\times$ \\
\bottomrule
\end{tabular}
\end{center}

Three repeats each. Power state is the largest single source of variation we measured in this TD,
and we did not record it during the main campaign, so that baseline is not strictly reproducible
as it stands. We also noticed, though we have not confirmed it, that the memory-bound case seems
less sensitive to CPU throttling than the compute-bound ones, which would be consistent with DRAM
speed not being governed by CPU power state. Confirming that would need a controlled comparison
with power state pinned, which we have not done yet.

\section{Cross-machine comparison}
\label{sec:cross-machine}
We collected data from Amani's and Yuba's machines independently and merged it into this report.
The raw files live in \code{data/} alongside \code{ultra9-285h}'s.

\subsection{Normalized \code{cpu\_task} throughput}
Amani and Yuba both benchmarked at sizes close to the sheet's own suggested values rather than the
scaled-up sizes we used on \code{ultra9-285h}. Amani's \code{cpu\_task} at size 200{,}000 still
measured 0.01\,s, the same sub-0.25\,s unreliability we flagged in Section~\ref{sec:campaign}, now
confirmed on a second machine. Dividing runtime by the $n^{1.5}$ cost shape lets us compare
directly across the different sizes each of us used:

\begin{center}
\begin{tabular}{lccc}
\toprule
Machine & Size & Runtime & $t / n^{1.5}$ \\
\midrule
Hichem (16c) & 10{,}000{,}000 & 2.120\,s & $6.70\times10^{-11}$ \\
Yuba (8c/16t)  & 5{,}000{,}000 & 0.780\,s & $6.98\times10^{-11}$ \\
Amani (8c)   & 1{,}000{,}000  & 0.09\,s  & $9.00\times10^{-11}$ \\
\bottomrule
\end{tabular}
\end{center}

Hichem's and Yuba's machines land within about 4\% of each other; Amani's machine is about 30 to
34\% slower per unit of work. Since \code{cpu\_task} is single-threaded, this isolates per-core
speed, which we think plausibly reflects the 268V being a lower-power (``V-series'') chip against
the 285H's and 11800H's higher-power designs, rather than a core-count effect. We have not
confirmed this against official specifications.

\subsection{Cache cliff, extended across our machines}
\begin{center}
\begin{tabular}{lccc>{\raggedright\arraybackslash}p{3.6cm}}
\toprule
Machine & L3 & Largest $N$ tested & Predicted working set & vs.\ L3 \\
\midrule
Hichem & 24\,MiB & 1500 & 54\,MB & 2.25$\times$ over: cliff observed \\
Amani  & 12\,MiB & 600  & 8.24\,MB & fits (69\% of L3): no cliff expected yet \\
Yuba   & 24\,MiB & 1000 & 24\,MB & at the boundary: cliff plausible \\
\bottomrule
\end{tabular}
\end{center}

Amani's predicted cliff, using her 12\,MiB L3, is around $N \approx 707$. We have not tested that
in her campaign, so we treat it as a prediction for a follow-up run rather than a confirmed result.
Yuba's only tested size ($N=1000$) sits almost exactly at his measured 24\,MiB L3 boundary, the
same boundary where we first saw the cliff on Hichem's machine in Finding 1. Yuba's
\code{matrix\_task} $N=1000$ mean is 0.947\,s against Hichem's 0.550\,s at the same $N$, which we
think is plausibly a mix of the general throughput gap above and sitting right at his cache
boundary rather than comfortably inside it: Hichem had 6.5\,MB of L3 headroom at $N=1000$, while
Yuba has essentially none. A run at $N=1500$ on Yuba's machine would confirm whether his cliff
lands where we predict.

\subsection{\code{memory\_task} across our machines}
This one was harder for us to compare cleanly. Amani's first repetition at each size ran
noticeably slower than her second and third (for example, size 512: 1.15\,s, then 0.70\,s,
0.66\,s), a first-run effect we also see more weakly in Hichem's own data. Using the full
three-run mean per the sheet's procedure, we get Hichem at about 0.0011\,s/MB, Yuba at about
0.00072\,s/MB, and Amani at about 0.00163\,s/MB (about 0.00133\,s/MB even discounting her outlier
first run). Amani's machine looks slower here too, but the effect is noisier than
\code{cpu\_task}'s, so we do not treat it as confirmed. It would need repeats with a stable,
recorded power state (Finding 2) to separate machine speed from run-to-run variance.

\section{Conclusion}
As a group, we found that the four workloads stress four different bottlenecks, confirmed by
measurement rather than assumed from the sheet's labels. \code{cpu\_task} is pure compute, scaling
almost exactly with its $n^{1.5}$ cost shape. \code{memory\_task} is allocation-bound, not
compute-bound: system time exceeds user time because the kernel is supplying and zeroing pages.
\code{matrix\_task} looks CPU-bound at small sizes but is not: once its working set exceeds L3,
doubling the input costs $8.19\times$ instead of the predicted $3.38\times$. \code{sort\_task} is
the mixed case, with the largest gap (10\%) between its internal timer and wall clock, which taught
us to trust the external GNU \code{time} figure over each program's self-reported number.

The single biggest lesson we take from this TD is not about any one workload: power state moved
every runtime by $\sim\!2\times$, dwarfing the differences between workloads at fixed input size.
Extending our comparison to three machines adds a second lesson: raw runtime is not comparable
across machines unless normalized by input size against the workload's known cost shape, or Amani's
faster-looking \code{cpu\_task} numbers would be misleading. We believe any later VM comparison will
be meaningless unless we record power state and machine specification alongside every measurement,
exactly the discipline this TD asked of us from the start.

\end{document}
